\documentclass[10pt,a4paper]{amsart}
\usepackage[margin=19mm]{geometry}
\usepackage{amsmath,amssymb,mathtools}
\usepackage[scr=rsfs,cal=euler]{mathalfa}
\usepackage{xfrac}
\usepackage[unicode,hidelinks,bookmarks=false]{hyperref}
\newcommand{\ie}{i.e.\ }
\newcommand{\cf}{cf.\ }
\newcommand{\resp}{respectively\ }
\newcommand{\Subsection}{\subsection}
\providecommand{\nobreakdash}{\nobreak\hbox{-}}
\setlength{\emergencystretch}{2em}
\multlinegap0pt
\usepackage{amssymb}
\usepackage{xcolor}
\newtheorem{theorem}{Theorem}
\newcommand{\mex}{\operatorname{mex}}
\title{Two-colored generalized Frobenius partitions and minimal-excludant sums over bipartitions}
\author{Rong Chen, Kang-Yu Wang,, Yi-ning Wang}
\date{Source: arXiv:2606.19696v1 (CC BY 4.0)}
\begin{document}
\maketitle

\noindent\emph{Source and licence.} Two-colored generalized Frobenius partitions and minimal-excludant sums over bipartitions. Version arXiv:2606.19696v1 (CC BY 4.0), distributed by arXiv under the Creative Commons Attribution 4.0 International licence, http://creativecommons.org/licenses/by/4.0/, as recorded in the arXiv licence metadata for this exact version and shown on the abstract page https://arxiv.org/abs/2606.19696v1. Full licence notice: \texttt{licenses/arxiv-2606.19696-CC-BY-4.0.txt}.\par\smallskip

\noindent\textit{Editorial scope.} Counted unit: the article's theorem thm.sigma.p2 (source0.tex lines 485--494). The article's complete original proof of it is delivered with it (source0.tex lines 496--669, one \texttt{proof} environment of 174 lines). No mathematics is written, repaired or re-derived: the statement and the proof are the article's own lines, in the article's own notation, and no retained line is substituted or re-flowed.  No further source-local context is needed: every cross-reference of every retained line resolves inside the two delivered spans.  Nothing was omitted from any retained span, so the package carries no omission record.  No retained line contains a citation, so the wrapper carries no bibliography and no bibliography entry.  No retained line contains a figure environment, an image-inclusion command or drawing code, so no image is delivered and the package declares no asset dependency.  Editorial formatting only, in addition: an AMS \texttt{amsart} preamble that restates the article's own declarations and macros used by the retained lines, namely the environments \texttt{theorem} on separate counters, together with the standard packages those lines need.  The retained environments print in this document as Theorem 1 in order of appearance, whereas the article prints the same items with the numbering of its own full document.  The complete native source of the article as distributed in the arXiv e-print for this exact version is bundled unchanged as \texttt{sources/203164/source0.tex}.
\begin{theorem}\label{thm.sigma.p2}
For all $n\geq 0$, we have
\[
2\sigma\mex_2(n)=\sum_{d\in\mathbb Z}p_2\bigl(n-d(d+1)\bigr)
\]
and
\[
2\sigma\mex_2(n)-E_2(n)=\sum_{d\in\mathbb Z}p_2(n-d^2).
\]
\end{theorem}

\begin{proof}
Let
\[
\delta(k)=(k,k-1,\ldots,2,1)
\]
be the staircase partition with largest part $k$, and put $\delta(0)=\emptyset$.  Then
\[
|\delta(k)|=1+2+\cdots+k=\frac{k(k+1)}2.
\]
For $N\in\mathbb Z$, let $\mathcal B(N)$ denote the set of bipartitions of $N$, with the convention that $\mathcal B(N)=\emptyset$ if $N<0$.  Thus $|\mathcal B(N)|=p_2(N)$.

We first prove the identity involving $\sigma\mex_2(n)$.  For $d\in\mathbb Z$, define
\[
\ell(d)=
\begin{cases}
d,&d\geq 0,\\
-d-1,&d<0.
\end{cases}
\]
Then
\[
d(d+1)=\ell(d)(\ell(d)+1)=2|\delta(\ell(d))|.
\]
For fixed $n$ and $d$, define
\[
\varphi_{n,d}:\mathcal B\bigl(n-d(d+1)\bigr)\longrightarrow \mathcal B(n)
\]
by
\[
\varphi_{n,d}(\mu_1,\mu_2)=
(\mu_1\cup\delta(\ell(d)),\mu_2\cup\delta(\ell(d))).
\]
This map adds one copy of each of the parts $1,2,\ldots,\ell(d)$ to both components.  Hence the image consists exactly of those bipartitions $\pi=(\pi_1,\pi_2)$ of $n$ for which each of the parts $1,2,\ldots,\ell(d)$ occurs in both $\pi_1$ and $\pi_2$.  Equivalently,
\[
\operatorname{Im}(\varphi_{n,d})
=A_{n,d}:=
\{\pi=(\pi_1,\pi_2)\vdash_2 n: \mex_2(\pi)>\ell(d)\}.
\]
The inverse map on $A_{n,d}$ is obtained by removing one copy of each of the parts $1,2,\ldots,\ell(d)$ from both components.  Thus $\varphi_{n,d}$ is a bijection from $\mathcal B(n-d(d+1))$ onto $A_{n,d}$, and therefore
\[
|A_{n,d}|=p_2\bigl(n-d(d+1)\bigr).
\]
Consequently,
\[
\sum_{d\in\mathbb Z}p_2\bigl(n-d(d+1)\bigr)
=
\sum_{d\in\mathbb Z}|A_{n,d}|.
\]
The right-hand side counts pairs $(d,\pi)$ such that $d\in\mathbb Z$, $\pi\vdash_2 n$, and $\pi\in A_{n,d}$.  We may count the same pairs by fixing $\pi$ first.  If $\mex_2(\pi)=m$, then $\pi\in A_{n,d}$ precisely when $\ell(d)<m$.  There are exactly $2m$ integers $d$ satisfying this condition, namely
\[
d=0,1,\ldots,m-1
\qquad\text{and}\qquad
 d=-1,-2,\ldots,-m.
\]
Thus each bipartition $\pi$ contributes $2\mex_2(\pi)$ to the sum.  Hence
\[
\sum_{d\in\mathbb Z}p_2\bigl(n-d(d+1)\bigr)
=
\sum_{\pi\vdash_2 n}2\mex_2(\pi)
=
2\sigma\mex_2(n).
\]

We now prove the identity involving $E_2(n)$.  The term $d=0$ in
\[
\sum_{d\in\mathbb Z}p_2(n-d^2)
\]
counts all bipartitions of $n$.  For $k\geq 1$, the term $d=k$ is represented by the map
\[
\psi^+_{n,k}:\mathcal B(n-k^2)\longrightarrow \mathcal B(n),
\]
where
\[
\psi^+_{n,k}(\mu_1,\mu_2)
=
(\mu_1\cup\delta(k),\mu_2\cup\delta(k-1)),
\]
and the term $d=-k$ is represented by the map
\[
\psi^-_{n,k}:\mathcal B(n-k^2)\longrightarrow \mathcal B(n),
\]
where
\[
\psi^-_{n,k}(\mu_1,\mu_2)
=
(\mu_1\cup\delta(k-1),\mu_2\cup\delta(k)).
\]
Both maps increase the weight by $k^2$, since
\[
|\delta(k)|+|\delta(k-1)|
=
\frac{k(k+1)}2+\frac{k(k-1)}2
=k^2.
\]
The image of $\psi^+_{n,k}$ is
\[
B^+_{n,k}
=
\{\pi=(\pi_1,\pi_2)\vdash_2 n:
\mex(\pi_1)>k,\ \mex(\pi_2)\geq k\},
\]
because one removes one copy of each of $1,2,\ldots,k$ from the first component and one copy of each of $1,2,\ldots,k-1$ from the second component to recover the preimage.  Similarly, the image of $\psi^-_{n,k}$ is
\[
B^-_{n,k}
=
\{\pi=(\pi_1,\pi_2)\vdash_2 n:
\mex(\pi_1)\geq k,\ \mex(\pi_2)>k\}.
\]
Thus $\psi^+_{n,k}$ and $\psi^-_{n,k}$ are bijections from $\mathcal B(n-k^2)$ onto $B^+_{n,k}$ and $B^-_{n,k}$, respectively.  It follows that
\[
\sum_{d\in\mathbb Z}p_2(n-d^2)
=
|\mathcal B(n)|+
\sum_{k\geq 1}|B^+_{n,k}|+
\sum_{k\geq 1}|B^-_{n,k}|.
\]
This identity is the point at which the bijections convert the $p_2$-sum into a weighted count of bipartitions of $n$: the left-hand side counts preimages in the various sets $\mathcal B(n-d^2)$, while the right-hand side counts their images in $\mathcal B(n)$, with multiplicity according to how many of the above image sets contain a fixed bipartition.

Fix a bipartition $\pi=(\pi_1,\pi_2)$ of $n$, and put
\[
m_1=\mex(\pi_1),
\qquad
m_2=\mex(\pi_2).
\]
The bipartition $\pi$ is counted once by the term $d=0$.  For $k\geq 1$, it is counted by $B^+_{n,k}$ precisely when
\[
m_1>k,
\qquad
m_2\geq k,
\]
and it is counted by $B^-_{n,k}$ precisely when
\[
m_1\geq k,
\qquad
m_2>k.
\]
Therefore the total multiplicity of $\pi$ in the above sum is
\[
1+\#\{k\geq 1:m_1>k,\ m_2\geq k\}
 +\#\{k\geq 1:m_1\geq k,\ m_2>k\}.
\]
The two cardinalities are
\[
\#\{k\geq 1:m_1>k,\ m_2\geq k\}=\min{m_1-1,m_2}
\]
and
\[
\#\{k\geq 1:m_1\geq k,\ m_2>k\}=\min{m_1,m_2-1}.
\]
Thus the total multiplicity is
\[
1+\min{m_1-1,m_2}+\min{m_1,m_2-1}.
\]
If $m_1=m_2=m$, this is $2m-1$.  If $m_1\neq m_2$, this is $2\min\{m_1,m_2\}$.  Since
\[
\mex_2(\pi)=\min\{m_1,m_2\},
\]
the multiplicity can be written uniformly as
\[
2\mex_2(\pi)-\mathbf{1}_{\{\mex(\pi_1)=\mex(\pi_2)\}},
\]
where $\mathbf{1}_E$ is $1$ if the condition $E$ holds and is $0$ otherwise.  Summing this multiplicity over all bipartitions $\pi$ of $n$ gives
\[
\sum_{d\in\mathbb Z}p_2(n-d^2)
=
\sum_{\pi\vdash_2 n}
\left(2\mex_2(\pi)-\mathbf{1}_{\{\mex(\pi_1)=\mex(\pi_2)\}}\right).
\]
By the definitions of $\sigma\mex_2(n)$ and $E_2(n)$, the last expression is
\[
2\sigma\mex_2(n)-E_2(n).
\]
This completes the proof.
\end{proof}

\end{document}
